\documentclass{article}
\usepackage[margin=25mm]{geometry}
\usepackage{amsmath,amssymb}
\usepackage{xurl}
\usepackage[hidelinks]{hyperref}
\setlength{\emergencystretch}{2em}
% Shared commands for notes in docs/paper-gaps/.

\DeclareUnicodeCharacter{2080}{\ifmmode{}_{0}\else\textsubscript{0}\fi}
\DeclareUnicodeCharacter{2081}{\ifmmode{}_{1}\else\textsubscript{1}\fi}
\DeclareUnicodeCharacter{2082}{\ifmmode{}_{2}\else\textsubscript{2}\fi}
\DeclareUnicodeCharacter{2099}{\ifmmode{}_{n}\else\textsubscript{n}\fi}

\newcommand{\Lean}{\textsc{Lean}}
\ExplSyntaxOn
\NewDocumentCommand{\leanid}{m}{
  \ifmmode
    \textsf{\detokenize{#1}}
  \else
    \group_begin:
    \urlstyle{sf}
    \tl_set:Nn \l_tmpa_tl {#1}
    \tl_replace_all:Nnn \l_tmpa_tl {\_} {_}
    \exp_args:NV \path \l_tmpa_tl
    \group_end:
  \fi
}
\ExplSyntaxOff
\providecommand{\tr}{\operatorname{tr}}
\newcommand{\tnleanIssueBase}{https://github.com/LionSR/TNLean/issues/}
\newcommand{\tnleanPrBase}{https://github.com/LionSR/TNLean/pull/}
\newcommand{\tnleanDocBase}{https://sirui-lu.com/TNLean/docs/}
\newcommand{\ghissue}[1]{\href{\tnleanIssueBase#1}{GitHub issue \##1}}
\newcommand{\ghpr}[1]{\href{\tnleanPrBase#1}{PR \##1}}
\newcommand{\doclink}[1]{\href{\tnleanDocBase#1.html}{\path{#1}}}

% Dirac notation and the identity operator, as in blueprint/src/macros/common.tex.
% \providecommand keeps note-local definitions authoritative.
\usepackage{dsfont}
\providecommand{\ket}[1]{|#1\rangle}
\providecommand{\bra}[1]{\langle #1|}
\providecommand{\braket}[2]{\langle #1 | #2 \rangle}
\providecommand{\ketbra}[2]{|#1\rangle\!\langle #2|}
\providecommand{\Id}{\mathds{1}}
\providecommand{\one}{\mathds{1}}
\providecommand{\C}{\mathbb{C}}
\providecommand{\MN}[1]{M_{#1}(\C)}
\providecommand{\MD}{M_D(\C)}

% Verdict marker: \gapnote{<kind>}{<status>}. Kind is the policy
% classification (clarification, local-correction, scope-restriction,
% unfaithful, false-source, open-gap); status is open, wip (elimination
% underway), resolved, or historical. Machine-read by the site index;
% prints a verdict line.
\newcommand{\gapnote}[2]{\par\noindent\textbf{Verdict:} #1 (#2).\par}

\title{The efficient approximate ground-space projector under (LD)}
\author{}
\date{2026-11-18}
\begin{document}
\maketitle
\gapnote{scope-restriction}{open}

\section{The source statement}

Theorem~5.9 (\texttt{t:aisgaplessefficient}) of Arad, Landau, Vazirani and
Vidick~\cite{Arad2016RigorousRG}, arXiv:1602.08828v2, file
\path{local_alg_arxiv_v2.tex}, lines 1424--1437, reads, with its proof at
lines 1439--1457:
\begin{quote}
Let $H$ be a local Hamiltonian satisfying Assumption (LD), parameters
$\eta_1\leq \eta$ and $\mu>0$, $\{1,\ldots,n\} = L\cup M\cup R$, where
$L=\{1,\ldots,i_1\}$, $M=\{i_1+1,\ldots,i_2\}$, and $R=\{i_2+1,\ldots,n\}$, a
partition of the $n$-qudit space, and $\varepsilon'_M$ an estimate for the
minimal energy $\varepsilon_M$ of the restriction of $H$ to $\mathcal H_M$
such that $|\varepsilon_M-\varepsilon'_M|\leq 10$. Then the procedure
\textsc{Generate}$(H, M, \varepsilon'_M,(\eta_1,\mu))$ described in Figure
\texttt{fig:generate} returns
\begin{itemize}
 \item MPO representations for a collection of $D^2$ operators
  $\{A_i\}_{i=1}^{D^2}$ acting on $\mathcal H_M$ and of bond dimension at most
  $e^{\tilde O(\log^3 n)}$ each such that there exists subspaces
  $\tilde T,\tilde T_-$ for which the conclusions of Theorem
  \texttt{t:aisgapless} are satisfied;
 \item An MPO for an operator $\tilde H_M$ such that
  $\|\tilde H_M\| = \tilde O(\log(n)/\mu)$ and the minimal energy
  $\tilde\varepsilon_M$ of $\tilde H_M$ restricted to $\tilde T_-$ satisfies
  $|\varepsilon_M-\tilde\varepsilon_M|<1/2$.
\end{itemize}
Moreover, \textsc{Generate}$(H, M, \varepsilon'_M, (\eta_1,\mu))$ runs in time
$e^{\tilde O(\log^3 n)}$.
\end{quote}
The quotation keeps the source's labelling of the sites by $1,\ldots,n$. The
running time is out of scope; only the existence content of the two items is
considered.

Throughout, $\varepsilon_0$ is the least eigenvalue of $H$, and $H_{[a,b]}$ is
the span of the eigenvectors of $H$ with eigenvalues in $[a,b]$.

The conclusions referred to are those of Theorem~4.8
(\texttt{t:aisgapless}, lines 1027--1035, proof 1037--1041): with $\mathcal H=\mathcal H_L\otimes
\mathcal H_M\otimes\mathcal H_R$, there are $D^2$ operators $A_i$ on
$\mathcal H_M$ and subspaces $\tilde T_-\subseteq\tilde T\subseteq\mathcal H$
such that
\begin{align}
 &H_{[\varepsilon_0,\varepsilon_0+\eta_1]}\ \text{is $.005$-close to}\ \tilde T,
 \qquad
 \tilde T_-\ \text{is $.005$-close to}\
 H_{[\varepsilon_0,\varepsilon_0+\eta_1-\mu/\log n]},
 \label{eq:alvv16_eld_agsp_windows}\\
 &D=e^{\tilde O(\frac{\log n}{\mu}\log^3 d)},\qquad D^{12}\Delta<10^{-5},
 \label{eq:alvv16_eld_agsp_count}
\end{align}
and every $S\subseteq\mathcal H_M$ that is $\delta$-viable for $\tilde T$ has
$\operatorname{span}\bigcup_iA_iS$ that is $\Delta/(1-\delta)^2$-viable for
$\tilde T_-$.

\section{A qualifier absent from the statement}

Theorem \texttt{t:aisgapless} assumes that $\mu>0$ is a constant and that
$\eta\geq\eta_1\geq2\mu/\log n$ (line 1027). Theorem~5.9 requires only
$\eta_1\leq\eta$ and $\mu>0$, while invoking the conclusions of Theorem
\texttt{t:aisgapless}. The source attaches the lower bound explicitly at that
comparable existence statement, and neither the statement of Theorem~5.9 nor
the description of \textsc{Generate} at line 1348 repeats it. Following the
convention that silence is not always an omission, the faithful reading of
Theorem~5.9 omits the lower bound. Whether the conclusions
(\ref{eq:alvv16_eld_agsp_windows}) still hold for $\eta_1<2\mu/\log n$ is an
open question for the formal statement: for $\eta_1<\mu/\log n$ the lower window
lies below $\varepsilon_0$, and the second closeness claim concerns the zero
space.

\section{What is formalized}

The formalization constructs, without any spectral gap, a soft-truncated
low-density improving family of middle matrix product operators
(\leanid{TNLean.MPS.AreaLaw.ALVV.exists_efficient_soft_low_density_family}) and
a calibration of it
(\leanid{TNLean.MPS.AreaLaw.ALVV.exists_calibrated_efficient_soft_low_density_family}).
The calibrated family gives the window structure of
(\ref{eq:alvv16_eld_agsp_windows}) with losses at most $1/200$, the bound
$D^{12}\Delta\leq10^{-5}$ in place of the strict inequality, and the viability
improvement, for a fixed buffer $0<w\leq1$ and band width at most $34$. The
source's buffer is $\mu/\log n$, which shrinks with the chain length; the
calibrated parameters then depend on $n$. These results are restricted versions
of the window and improvement clauses, not the source theorem.

The following clauses of Theorem~5.9 are not formalized:
\begin{enumerate}
 \item the member bond dimension $e^{\tilde O(\log^3 n)}$ when the buffer is
  $\mu/\log n$;
 \item the count $D=e^{\tilde O(\frac{\log n}{\mu}\log^3 d)}$ at that buffer;
 \item the norm bound $\|\tilde H_M\|=\tilde O(\log(n)/\mu)$;
 \item the energy estimate $|\varepsilon_M-\tilde\varepsilon_M|<1/2$, the
  uniform energy estimate recorded as open in
  \href{https://sirui-lu.com/TNLean/paper-gaps/alvv16_soft_regional_residuals.pdf}{the soft-residual note}.
\end{enumerate}
The source proof (lines 1439--1457) chooses
\begin{align}
 \ell=\Theta\Big(\frac{\log n}{\mu}\log\frac{\log n}{\mu}\Big),\qquad
 k=\ell^2,\qquad t=\Theta(\ell),\qquad t'=4,
 \label{eq:alvv16_eld_agsp_parameters}
\end{align}
and the fourth-order truncation $t'=4$ is not an instance of the formalized
calibration, as recorded in
\href{https://sirui-lu.com/TNLean/paper-gaps/alvv16_soft_regional_residuals.pdf}{the soft-residual note}. Here the soft-truncation
parameters are written $t,t'$, as in the source.

\section{Reading of the asymptotic bounds}

A formal statement must fix the meaning of $\tilde O$. The bounds are read with
constants depending only on the local dimension $d$, the energy parameter $\eta$
of Assumption (LD) and $\mu$, uniformly in $n$, the partition and the estimate
$\varepsilon'_M$; $\tilde O(f)$ allows a factor polylogarithmic in $f$. With
this reading the first item asks for a bound
\begin{align}
 \log(\text{bond of }A_i)\leq C\log^3 n\,(\log\log n)^{C}
 \label{eq:alvv16_eld_agsp_bond}
\end{align}
for all $n\geq3$, with $C$ depending only on $d$, $\eta$ and $\mu$.

\section{What the printed proof supports}

The following are findings about the printed proof, not refutations of
Theorem~5.9.

\paragraph{The bond bound.} The formalized calibration chooses one natural
number $u$ from the local dimension and the effective gap $g=w^2/2$ of the
buffer $w$, and sets
\begin{align}
 \ell=u^8,\qquad m=u^{10},\qquad k=u^{16},
 \label{eq:alvv16_eld_agsp_calibration}
\end{align}
where $m$ is the soft-truncation order and $k$ the filter degree. The member
bond is $\bigl(2(m+1)\,o^2q^{o}\bigr)^{2k}$, where the cluster order $o$ grows
like $\log N$ in the number $N$ of bonds, so its logarithm is of order
$u^{16}\log N$. At the source's buffer $w=\mu/(3\log n)$ the number $u$ grows
like a power of $\log n/\mu$, and the logarithm of the bond is a power of
$\log n$ larger than three; the bond is quasi-polynomial, but not
$e^{\tilde O(\log^3n)}$. The printed bound needs the parameters
(\ref{eq:alvv16_eld_agsp_parameters}), in particular the fourth-order truncation
$t'=4$, and
\leanid{TNLean.MPS.AreaLaw.ALVV.not_soft_fourth_order_residual_budget} shows that
the fourth order fails the residual budget of the formalized calibration.
Recovering the printed bound therefore needs a sharper residual estimate than
the one formalized.

\paragraph{The energy estimate.} The proof derives
$|\varepsilon_M-\tilde\varepsilon_M|<1/2$ in one sentence, at line~1445:
``The claim on the ground state energies of $H_M$ and $\tilde H_M$ follows
directly from Lemma~\texttt{lem:exp-approx} as well.'' The first estimate of that
lemma, equation \texttt{eq:exp-approx-0} at lines 1195--1201, is shown false in
its uniform reading in
\href{https://sirui-lu.com/TNLean/paper-gaps/alvv16_soft_truncation_energy_estimate.pdf}{the
soft-truncation energy note}, for every fixed admissible order and scale. Only a
corrected comparison for one interval is formalized. The energy clause therefore
needs its own argument; the printed one does not establish it.

\section{Verdict}

Scope restriction, open. The formalized families prove restricted versions of
the window and improvement clauses for a fixed buffer. Theorem~5.9 itself, with
its shrinking buffer, its bond and norm bounds and its energy estimate, is not
formalized. The lower bound $\eta_1\geq2\mu/\log n$ is not added as a
hypothesis. The analysis of the printed proof in the preceding section records
why the bond bound and the energy estimate are not reached by the formalized
families; it does not decide whether they hold.

\bibliographystyle{plain}
\bibliography{references}

\end{document}
